% ============================================================
% Come si usano, argomento 16: vettori aleatori e vettori gaussiani
% ============================================================
\argomento{Vettori aleatori e vettori gaussiani}{Formule a p.~\pageref*{f:vettore}}[Vettori aleatori e gaussiani]

\begin{scelta}{Come scegliere la formula}
\renewcommand{\arraystretch}{1.25}
\begin{tabular}{@{}P{0.55\linewidth}P{0.39\linewidth}@{}}
  \toprule
  \textbf{Se il testo\ldots} & \textbf{si usa\ldots} \\
  \midrule
  descrive molte variabili insieme (campioni, antenne, misure) & vettore delle medie e matrici \vv{vettmedie}
    \vv{matrcorr} \vv{matrcov} \vv{covcorr} \\
  chiede la varianza di una combinazione lineare delle variabili & la forma quadratica \vv{nonnegativa} \\
  dice che le variabili sono congiuntamente gaussiane & la densità \vv{vettgauss}; con due variabili
    \vv{bivariata} \\
  chiede se variabili gaussiane incorrelate sono indipendenti & sì, se congiuntamente gaussiane \vv{gaussincorr} \\
  chiede la distribuzione di una variabile sapendo l'altra & la condizionata \vv{gausscond} \\
  applica una trasformazione lineare (somme, differenze, filtri, medie) & \vv{trasflineare}, \vv{mediagauss} \\
  chiede come generare variabili correlate & \vv{generacorr} \\
  \bottomrule
\end{tabular}

Il controllo delle dimensioni salva da molti errori: in $\matr{A}\matr{C}_{\vett{X}}\matr{A}^T$ le dimensioni interne
devono coincidere, $(m \times n)(n \times n)(n \times m) = m \times m$.
\end{scelta}

\scheda{vettore}{
  \usoquando{Quando le variabili sono molte e conviene trattarle insieme: i campioni di un segnale, le tensioni su più
    antenne, più misure della stessa grandezza.}
  \usoesempio{Due misure $X_1$, $X_2$ formano $\vett{X} = [X_1, X_2]^T$; la sua densità è la congiunta del capitolo 14.}
}

\scheda{vettmedie}{
  \usoquando{Per riassumere dove si trova il vettore: una media per ogni componente.}
  \usoesempio{Densità $x_1 + x_2$ su $[0, 1]^2$: $\vett{\mu}_{\vett{X}} = \left[\frac{7}{12}, \frac{7}{12}\right]^T$.}
}

\scheda{matrcorr}{
  \usoquando{Raccoglie in una tabella tutti i valori $E[X_i X_j]$; serve per arrivare alla matrice di covarianza.}
  \usoesempio{Densità $x_1 + x_2$: $E[X_i^2] = \frac{5}{12}$, $E[X_1 X_2] = \frac{1}{3}$, quindi
    $\matr{R}_{\vett{X}} = \frac{1}{144}\begin{bmatrix} 60 & 48 \\ 48 & 60 \end{bmatrix}$.}
}

\scheda{matrcov}{
  \usoquando{È la seconda metà della caratterizzazione sintetica: varianze sulla diagonale, covarianze fuori.}
  \usoattenzione{È sempre simmetrica e ha la diagonale positiva: se non lo è, c'è un errore.}
  \usoesempio{Due variabili con $\sigma_1 = 2$, $\sigma_2 = 3$ e $\rho = 0{,}5$: $\matr{C}_{\vett{X}} =
    \begin{bmatrix} 4 & 3 \\ 3 & 9 \end{bmatrix}$.}
}

\scheda{covcorr}{
  \usoquando{Si conoscono le medie e la matrice di correlazione (o è più facile calcolare quella): si toglie il
    prodotto colonna per riga delle medie.}
  \usoesempio{Densità $x_1 + x_2$: $\vett{\mu}_{\vett{X}} \vett{\mu}_{\vett{X}}^T = \frac{1}{144}\begin{bmatrix} 49 &
    49 \\ 49 & 49 \end{bmatrix}$, quindi $\matr{C}_{\vett{X}} = \frac{1}{144}\begin{bmatrix} 11 & -1 \\ -1 & 11
    \end{bmatrix}$.}
}

\scheda{nonnegativa}{
  \usoquando{Si vuole la varianza di una combinazione lineare $\vett{q}^T \vett{X}$; oppure si deve controllare che una
    matrice possa essere di covarianza.}
  \usocome{Si calcola $\vett{q}^T \matr{C}_{\vett{X}} \vett{q}$ riga per colonna.}
  \usoesempio{Densità $x_1 + x_2$, $\vett{q} = [1, 1]^T$: $\mathrm{VAR}(X_1 + X_2) = \frac{11 - 1 - 1 + 11}{144} =
    \frac{20}{144}$. Una matrice $\begin{bmatrix} 1 & 2 \\ 2 & 1 \end{bmatrix}$ non può essere di covarianza: con
    $\vett{q} = [1, -1]^T$ darebbe una varianza di $-2$.}
}

\scheda{vettgauss}{
  \usoquando{Le variabili sono congiuntamente gaussiane: rumori, misure disturbate, campioni di un processo
    gaussiano.}
  \usocome{Bastano $\vett{\mu}_{\vett{X}}$ e $\matr{C}_{\vett{X}}$: si calcolano il determinante e l'inversa di
    $\matr{C}_{\vett{X}}$ e si scrive la forma quadratica all'esponente.}
  \usoesempio{Con $n = 1$ la matrice è il numero $\sigma_X^2$, e la formula ridà la gaussiana
    $\frac{1}{\sqrt{2\pi\sigma_X^2}}\, e^{-\frac{(x - \mu_X)^2}{2\sigma_X^2}}$.}
}

\scheda{bivariata}{
  \usoquando{Due variabili congiuntamente gaussiane di cui si conoscono medie, varianze e coefficiente di
    correlazione.}
  \usoattenzione{Serve $\abs{\rho} < 1$: con $\rho = \pm 1$ le variabili sono legate da una retta e la densità nel
    piano non esiste.}
  \usoesempio{$\sigma_1 = \sigma_2 = 1$, medie nulle, $\rho = 0{,}5$: $\det \matr{C}_{\vett{X}} = 0{,}75$ e
    $f(0, 0) = \frac{1}{2\pi\sqrt{0{,}75}} \simeq 0{,}184$.}
}

\scheda{cerchi}{
  \usoquando{Per la probabilità che un punto gaussiano cada vicino al centro: la regola delle tre sigma in due
    dimensioni.}
  \usocome{Con $\rho = 0$ e varianze unitarie si integra in coordinate polari (\vv{polari}).}
  \usoesempio{Entro raggio 2 dal centro: $1 - e^{-2} \simeq 0{,}865$, meno del $0{,}954$ della regola delle due sigma in
    una dimensione.}
}

\scheda{gaussincorr}{
  \usoquando{Si devono dimostrare o usare indipendenze fra variabili gaussiane: per le congiuntamente gaussiane basta
    controllare che le covarianze siano nulle.}
  \usoattenzione{Devono essere \emph{congiuntamente} gaussiane. $X_1 \sim \mathcal{N}(0, 1)$ e $X_2 = S X_1$, con $S = \pm
    1$ a caso: entrambe gaussiane e incorrelate, ma $\abs{X_2} = \abs{X_1}$ sempre.}
  \usoesempio{Somma e differenza di due rumori gaussiani indipendenti con la stessa varianza hanno covarianza nulla:
    sono indipendenti.}
}

\scheda{gausscond}{
  \usoquando{Si osserva una delle due variabili e si vuole dire qualcosa sull'altra: la media condizionata è la
    previsione, la varianza condizionata l'incertezza che resta.}
  \usocome{Nel caso generale si standardizzano le due variabili, si usa la formula e si torna indietro.}
  \usoesempio{$\rho = 0{,}8$, osservato $X_1 = 1$: $X_2 \sim \mathcal{N}(0{,}8;\ 0{,}36)$, cioè previsione $0{,}8$ con
    deviazione standard $0{,}6$ invece di 1.}
}

\scheda{proprietagauss}{
  \usoquando{Per sapere che cosa si può dare per scontato con le gaussiane: tutto ciò che si ottiene con operazioni
    lineari, sottoinsiemi e condizionamenti resta gaussiano, e basta calcolarne medie e covarianze.}
  \usoesempio{Il segnale ricevuto $Y = A + N$ con rumore gaussiano è gaussiano; anche la media di più campioni lo è.}
}

\scheda{trasflineare}{
  \usoquando{Il vettore passa attraverso un sistema lineare: somme, differenze, medie, filtri, cambi di coordinate.}
  \usocome{Media: $\matr{A}\vett{\mu}_{\vett{X}} + \vett{b}$. Covarianza: $\matr{A}\matr{C}_{\vett{X}}\matr{A}^T$ (la
    costante $\vett{b}$ non conta). Se $\vett{X}$ è gaussiano, l'uscita è gaussiana con quei parametri.}
  \usoesempio{Due rumori indipendenti $\mathcal{N}(0, \sigma^2)$ e $\matr{A} = \begin{bmatrix} 1 & 1 \\ 1 & -1
    \end{bmatrix}$: $\matr{C}_{\vett{Z}} = \sigma^2 \matr{A}\matr{A}^T = \begin{bmatrix} 2\sigma^2 & 0 \\ 0 & 2\sigma^2
    \end{bmatrix}$, somma e differenza incorrelate.}
}

\scheda{mediagauss}{
  \usoquando{Si mediano $n$ misure indipendenti disturbate da rumore gaussiano.}
  \usoesempio{Con 100 misure la deviazione standard del rumore diventa $\frac{\sigma}{10}$: l'errore tipico è dieci
    volte più piccolo.}
}

\scheda{generacorr}{
  \usoquando{Si devono simulare due variabili gaussiane con una correlazione assegnata, partendo da numeri gaussiani
    indipendenti.}
  \usoesempio{$\rho = 0{,}6$: $X_1 = W_1$ e $X_2 = 0{,}6\, W_1 + 0{,}8\, W_2$, perché $\sqrt{1 - 0{,}36} = 0{,}8$.}
}
